\section{Глава~\ref{sec:muscles}. Выход на мышцы}

Устройство выхода --- двигательные единицы, запас надёжности синапса, включение по размеру, пара мышц, петля растяжения и генераторы ритма --- измерено прямо; условие устойчивости петли с задержкой --- теорема; внутренняя модель тела --- хорошо обоснованная гипотеза; числа рисунков --- расчёт по модели книги.

{\footnotesize
\setlength{\tabcolsep}{3pt}\renewcommand{\arraystretch}{1.15}
\begin{longtable}{>{\raggedright}p{0.5cm}>{\raggedright}p{4.0cm}>{\raggedright}p{5.6cm}>{\raggedright}p{2.3cm}>{\raggedright\arraybackslash}p{1.7cm}}
\toprule
№ & Утверждение & Опыт и что измерено & Источник & Статус \\
\midrule\endfirsthead
\toprule
№ & Утверждение & Опыт и что измерено & Источник & Статус \\
\midrule\endhead
1 & Двигательная единица: один \nt{АЦЕТ}-нейрон управляет группой волокон, от нескольких сотен до пары тысяч; в мышцах тонкой подстройки единицы мельче & Мышцы человека, подсчёт двигательных аксонов и мышечных волокон: в среднем около $340$ волокон на нейрон в межкостной мышце кисти, около $410$ в плечелучевой мышце, около $1930$ в медиальной икроножной. Точного числа для самых мелких единиц глава не приводит. & \cite{Feinstein1955mu} & измерено \\
2 & Синапс на мышце выбрасывает в несколько раз больше медиатора, чем нужно для срабатывания волокна & Обзор измерений на нервно-мышечных синапсах: запас надёжности (отношение выброшенного медиатора к пороговому) у взрослых млекопитающих обычно $3$--$5$; у человека запас держится больше на складках мембраны получателя, чем на величине выброса. & \cite{WoodSlater2001} & измерено \\
3 & Подёргивание~\eqref{eq:twitch} нарастает за время $T_c$ порядка $50\ms$ & Одиночные двигательные единицы кисти человека при произвольном усилии, сила усреднена по импульсам единицы: время нарастания $30$--$100\ms$, у $80\,\%$ единиц меньше $70\ms$. Функция $(t/T_c)e^{1-t/T_c}$ --- аппроксимация из модели пула единиц. & \cite{MilnerBrown1973rec, Fuglevand1993} & измерено \\
4 & Сумма подёргиваний~\eqref{eq:forcesum}: при $5$, $15$ и $40$ имп/с сила $0.2$--$1.1F_1$, $1.8$--$2.2F_1$ и около $5.4F_1$ (рис.~\ref{fig:tetanus}) & Расчёт по \texttt{models/\allowbreak muscle\_\allowbreak models.py}, $T_c=50\ms$: $0.21$--$1.10$, $1.76$--$2.19$ и $5.28$--$5.49\,F_1$ на последних $0.3\,\text{с}$. Линейное сложение --- упрощение: насыщение настоящей мышцы в модели нет. & \texttt{models/\allowbreak muscle\_\allowbreak models.py} & модель \\
5 & Единицы включаются по размеру; самые слабые в сотню раз слабее самых сильных & Вентральные корешки кошки: при нарастающем рефлекторном входе первыми разряжаются нейроны с мелкими импульсами в корешке, то есть мелкие. Межкостная мышца кисти человека: сила подёргивания единиц от $0.1$ до $10$~г и растёт почти линейно с силой, при которой единица включается. & \cite{Henneman1957, MilnerBrown1973rec} & измерено \\
6 & Шаг силы пропорционален силе, $\Delta F/F\approx q-1$~\eqref{eq:sizeprinciple}: плавающая точка & Выведено в главе из геометрической прогрессии сил. Согласуется с опытом: при больших усилиях на то же приращение силы включается гораздо меньше новых единиц. & вывод в главе; \cite{MilnerBrown1973rec} & теория \\
7 & Разброс силы растёт пропорционально самой силе & Произвольное изометрическое усилие мышцы большого пальца: стандартное отклонение силы линейно растёт со средней силой; при том же усилии, вызванном электрической стимуляцией, разброс постоянен. Модель пула: линейный рост даёт включение единиц по порядку силы. & \cite{Jones2002sdn} & измерено \\
8 & Плавность движений --- следствие шума, зависящего от сигнала & Модель: траектория, минимизирующая разброс конечной точки при таком шуме, воспроизводит измеренные траектории глаза и руки. Прямого опыта, отличающего эту причину от других, нет. & \cite{HarrisWolpert1998} & гипотеза \\
9 & Разность сил пары мышц задаёт момент, сумма --- жёсткость~\eqref{eq:antagonists} & Теория управления импедансом совместным напряжением. Рука человека: жёсткость каждого сустава, измеренная малыми толчками, примерно пропорциональна моменту в нём; разные соотношения совместного напряжения меняют форму и направление эллипса жёсткости кисти. & \cite{Hogan1984, GomiOsu1998stiff} & измерено \\
10 & Датчик растяжения возбуждает свой \nt{АЦЕТ}-нейрон и через тормозящего посредника выключает противника (рис.~\ref{fig:antagonists}) & Спинной мозг кошки: одиночный посредник, возбуждаемый волокнами датчиков растяжения, даёт моносинаптические тормозные потенциалы в \nt{АЦЕТ}-нейронах мышцы-противника, синаптическая задержка $0.28$--$0.42\ms$. Медиатор посредника (глицин) в этом опыте не определялся. & \cite{Jankowska1972Ia} & измерено \\
11 & Задержка петли: спинномозговой $25$--$30\ms$, через кору в полтора--три раза больше, через глаз $100$--$200\ms$ & Рука человека, толчок по суставу: короткий ответ мышц через $20$--$45\ms$, длинный $45$--$105\ms$, произвольный $120$--$180\ms$. Смещение видимого положения пальца во время движения: поправка через $160\ms$ в среднем. & \cite{Pruszynski2008rapid, SaundersKnill2003vis} & измерено \\
12 & $\dd x/\dd t=-Gx(t-d)$ устойчиво при $Gd<\pi/2$~\eqref{eq:delaystable}; на границе частота $1/(4d)$ & Теорема о корнях уравнения с запаздыванием. Проверка по \texttt{models/\allowbreak muscle\_\allowbreak models.py}, $d=30\ms$: к $3\,\text{с}$ амплитуда $3\cdot10^{-6}$ при $G=40$, $0.13$ при $G=50$, $38$ при $G=55$ в секунду (граница $52$). & \cite{Hayes1950} & теория \\
13 & Физиологическая дрожь $8$--$12$~Гц; петля растяжения --- один из её источников & Обзор измерений: мелкая дрожь в диапазоне около $10$~Гц есть у здоровых людей; её происхождение смешанное --- центральные колебания, свойства разрядов единиц, механический резонанс и резонанс рефлекторной петли, и в разных условиях преобладает разное. & \cite{McAuleyMarsden2000} & гипотеза \\
14 & Мозг предсказывает результат команды по внутренней модели тела & Человек держит груз щипком и двигает рукой: при инерционной, вязкой и упругой нагрузке сила сжатия меняется одновременно с нагрузкой, а не с запаздыванием петли, то есть нагрузка предсказана. Обзор сводит такие опыты; прямого наблюдения самой модели в нейронах нет. & \cite{FlanaganWing1997grip, WolpertGhahramani2000} & гипотеза \\
15 & Ритм ходьбы, дыхания, жевания порождают местные сети при постоянном входе & Обзор опытов: изолированная нервная система (спинной мозг, ганглии) выдаёт ритмический двигательный узор без ритмического входа и без обратной связи от мышц. & \cite{MarderBucher2001} & измерено \\
16 & Взаимное торможение с усталостью~\eqref{eq:halfcentre} даёт ритм с периодом $0.84\,\text{с}\approx2\tau_a$ (рис.~\ref{fig:halfcentre}) & Теорема: сеть со взаимным торможением и адаптацией без устойчивого состояния при постоянном входе обязательно колеблется. Расчёт по \texttt{models/\allowbreak muscle\_\allowbreak models.py} ($I=1$, $\beta=2$, $b=2$, $\tau=20\ms$, $\tau_a=0.4\,\text{с}$): период $0.84\,\text{с}$. Что настоящие генераторы устроены именно так, не показано. & \cite{Matsuoka1985osc} & модель \\
\bottomrule
\end{longtable}}
